% ===== tr-06 块2：第 6 章 pp.17--21（原文页 17--21）=====
% [边界] p.16/17 跨页句 “If one replaces ω in (6.6.1) by Im Ω …” 已在 tr06d 末整句译完；本文件从 p.17 该句之后起译（页首两行为该句的收尾，不再重复）。
% [原文 p.17]
这些坐标在 SYZ 猜想中起关键作用------该猜想由 Gross 重新表述并精炼（见 [24, 猜想 6.6.6]）。为更详细地解释 SYZ 镜像对称，我们必须解释复模空间与 Kähler 模空间的相关概念。由于镜像对称预期发生在“大复结构极限”附近，我们需要在这一框架中定义大复结构极限的概念。

给定带光滑反典范除子 $D$ 的次数为 $d$ 的 del Pezzo 曲面 $Y$，我们在 $X=Y\setminus D$ 上固定精确 Tian--Yau 度量 $\omega_{TY}$。这对应于 $Y$ 上一个单调辛结构的选择。因此，镜像的有理椭圆曲面应当容许一个特殊的复结构。为描述这样的有理椭圆曲面，我们首先回顾 Gross--Hacking--Keel [26] 关于 Looijenga 对的 Torelli 定理，因为它在当前情形下适用。设 $\check Y$ 是带有 $I_d$ 纤维 $\check D$ 的有理椭圆曲面，令 $\check X=\check Y\setminus\check D$。则配对 $(\check Y,\check D)$ 的周期由
\begin{align}
\varphi_{(\check Y,\check D)}&:H_2(\check X)/\operatorname{Im}H^1(\check D)\to\C^*, \notag\\
\gamma&\mapsto\exp 2\pi i\left(\frac{\int_\gamma\check\Omega}{\int_{[C]}\check\Omega}\right). \tag{6.6.2}
\end{align}
给出。
% [Tutor Note: 周期映射的定义域印作 $H_2(\check X)/\operatorname{Im}H^1(\check D)$，即同调群模去上同调群在映射下的像------下标 $2$ 与上标 $1$ 并不一致（页 17 可逐字核对），原书如此；照录不改。另外这里 $(6.6.2)$ 是本章第 4 个编号式子中的第 2 个，全章共 4 个编号式（(6.6.1)--(6.6.4)）。]

由 (6.3.1)，我们有 $H_2(\check X)/\operatorname{Im}H^1(\check D)\cong\check D^{\perp}$。这里 $\check D^{\perp}$ 是 $\operatorname{Pic}(\check Y)$ 中与 $\check D$ 的每个分支配对为零的子群。于是 $\varphi_{(\check Y,\check D)}\in\operatorname{Hom}(\check D^{\perp},\C^*)$。由 [19, 命题 3.12] 中的留数计算可知：经由等同 $\operatorname{Pic}^0(\check D)\cong\C^*$，它与 $\check D^{\perp}$ 中的线丛限制到 $\check D$ 所得相同。有了上述理解，我们可以描述配对 $(\check Y,\check D)$ 的模空间：

\begin{theorem}{6.6.1}[{Gross--Hacking--Keel, [26]}]
设 $(\check Y',\check D')$ 与配对 $(\check Y,\check D)$ 处于同一形变族。我们通过平行移动把 $\check D^{\perp}$ 与 $\check D'^{\perp}$ 等同。则 $\varphi_{(\check Y',\check D')}=\varphi_{(\check Y,\check D)}$ 蕴含配对同构 $(\check Y',\check D)\cong(\check Y,\check D)$。此外，给定任意同态 $\varphi\in\operatorname{Hom}(\check D^{\perp},\C^*)$，存在同一形变族中周期为 $\varphi$ 的配对 $(\check Y',\check D')$。
\end{theorem}
% [Tutor Note: 定理中“蕴含配对同构 $(\check Y',\check D)\cong(\check Y,\check D)$”原文如此（页 17 页图可辨，第二个配对缺了 $D$ 的撇号）；按命题本意应为 $(\check Y',\check D')\cong(\check Y,\check D)$。译文照录原样，此注给出正确形式。]

这一结果意味着：带 $I_d$ 纤维的有理椭圆曲面的形变族中存在一个特殊配对 $(\check Y_e,\check D_e)$，其周期为平凡。我们记 $\check X_e=\check Y_e\setminus\check D_e$。我们断言，这个特殊的复结构正是与精确 Tian--Yau 度量的选择互为镜像的那个。

与 $(\check Y,\check D)$ 的情形类似，由 [53, 定理 6.4] 以及 [19, 命题 3.12] 的类似论证，配对 $(Y,D)$ 的周期在其形变族中决定了该配对。以 $\mathcal{M}_{cpx}$ 记带光滑反典范除子的 del Pezzo 曲面的带标号配对的模空间。存在自然的全纯映射 $\mathcal{M}_{cpx}\to\C$，由取光滑反典范除子的 $j$-不变量给出，而后者由 $\operatorname{Im}H^1(D_0)\to H_2(X_0)\cong\Z^2$ 的周期决定。\footnote{严格说来，我们需要考虑带标号配对连同 SYZ 纤维类以及使 SYZ 纤维相位为零的一个体积形式的模空间。}
% [原文 p.18]
这里 $(Y_0,D_0)$ 是参考配对，$X_0=Y_0\setminus D_0$。纤维于是双全纯于 $\operatorname{Hom}(D_0^{\perp},D_0)$，由 $H_2(X_0)$ 中其余的周期决定。注意这里 $D_0$ 是光滑椭圆曲线，因此有自然同构 $\operatorname{Pic}^0(D_0)\cong D_0$。

为 SYZ 镜像对称，我们将选取一个本原 2-圈 $[\check L]\in H^2(\check X_e,\Z)$，它可以由包含在 $\check D$ 的邻域中的一个 2-圈代表。我们定义复化 Kähler 模空间为
\[
\check{\mathcal{M}}_{Kah}=\{\mathbb{B}+i[\omega]\in H^2(\check X_e,\C)\mid[\check\omega]\cdot[\check L]=0,\ [\check\omega]\in\mathcal{K}_{CY}\}/H^2(\check X_e,\Z).
\]
% [Tutor Note: 式中 $\mathbb{B}$ 是 B-场的记号（黑板体 B，页 18 字体可辨），$[\check\omega]\in\mathcal{K}_{CY}$ 中的 $K$ 是花体 $\mathcal{K}_{CY}$（与 §5 中带 check 帽的 $\check K_{CY}$ 字体不同，原书两处并不统一）；照录。]

\begin{remark}{6.6.2}
不同 $[\check L]$ 选取给出的复化 Kähler 模空间 $\check{\mathcal{M}}_{Kah}$ 自然等同，并且与 $\mathcal{M}_{cpx}$ 上的单值作用互为镜像。
\end{remark}

我们仍需为 $\check{\mathcal{M}}_{cpx}$ 定义“大复结构极限”的概念，镜像对称预期发生在那里。设 $(Y_i,D_i)$ 是一列带反典范除子 $D_i$ 的 del Pezzo 曲面，使得
\begin{enumerate}
\item $Y_i$ 收敛到光滑 del Pezzo 曲面 $Y_\infty$，且
\item $D_i$ 收敛到 $Y_\infty$ 中的不可约结点反典范除子 $D_\infty$。
\end{enumerate}
在 [11] 中我们观察到：与 $D_i$ 的消没圈相对应的 $Y_i\setminus D_i$ 的特殊拉格朗日纤维化在无穷远处坍缩。这正是从 SYZ 观点看大复结构极限的特征 [44]；我们将说这样一列配对收敛到大复结构极限。

\begin{question}{6.6.3}
$X_i$ 的 SYZ 纤维化是否在 $i\to\infty$ 时整体坍缩？注意这里大复结构极限的定义并不是通常的模空间最深的边界分层。
\end{question}

有了上述几何框架，我们可以精确陈述 del Pezzo 曲面与有理椭圆曲面之间 SYZ 镜像对称的概念。

\begin{theorem}{6.6.4}[{C.-Jacob-L., [11]}]
充分接近大复结构极限时，存在镜像映射 $\Psi:\mathcal{M}_{cpx}\to\check{\mathcal{M}}_{Kah}$，使得 $X_q=Y_q\setminus D_q$ 与 $\check X_{\Psi(q)}=\check Y_{\Psi(q)}\setminus\check D_{\Psi(q)}$ 上的特殊拉格朗日纤维化在如下意义下互为对偶纤维化：
\begin{enumerate}
\item 设 $B_q$、$B_{\Psi(q)}$ 分别是 $X_q$、$\check X_{\Psi(q)}$ 中 SYZ 纤维化的底。存在微分同胚 $\Psi_q:B_q\to B_{\Psi(q)}$，使得 $\Psi$ 交换两个 SYZ 纤维化的复仿射结构与辛仿射结构。
\item 两个纤维的大小互为倒数。
\end{enumerate}
\end{theorem}

镜像映射 $\Psi$ 的选取使得定理 6.2.3 中 $X_q$ 的超 Kähler 旋转与定理 6.4.4 中 $\check X_{\Psi(q)}$ 的超 Kähler 旋转双全纯，但全纯体积形式相差一个相位 $\pi/2$，并被缩放以照顾纤维大小的关系。该双全纯可通过比较有理椭圆曲面相对于 $I_d$-纤维的周期并应用定理 6.6.1 来实现。总之，镜像配对由
% [原文 p.19]
两个超 Kähler 旋转相连接，它们匹配复结构但不匹配 Kähler 结构。

下面作者们提出一些将来会探讨的问题。第一个问题是非紧 Calabi--Yau 曲面的 SYZ 镜像对称与作为 Maslov 指标二圆盘加权计数的镜像超势之间的相容性。注意后者依赖于紧化。

\begin{question}{6.6.5}
给定 $\check X_e$ 上一个（复化的）Kähler 形式 $\check\omega_{\Psi(q)}$，镜像超势能否恢复 del Pezzo 曲面配对 $(Y_q,D_q)$？
\end{question}

定理 6.6.4 中的 SYZ 镜像对称相当于关掉 del Pezzo 曲面的 Kähler 模空间与有理曲面配对的复模空间。终极的 SYZ 猜想应同时包含 A-模型与 B-模型。

\begin{question}{6.6.6}
既然存在非恰当的 Tian--Yau 度量，定理 6.6.4 中的 SYZ 镜像对称能否推广到同时包含 del Pezzo 曲面的 Kähler 模空间与有理椭圆曲面的复模空间？
\end{question}

人们也可以反向地追问镜像对称不同方面的相容性。

\begin{question}{6.6.7}
若上一问题的答案是肯定的，则 SYZ 镜像是否与已知的、对应于 del Pezzo 曲面非单调辛形式的超势一致？
\end{question}

在本节末尾，我们将解释从特殊拉格朗日纤维化得到 del Pezzo 曲面超势的方法。考虑一列收敛到大复结构极限的配对 $(Y_i,D_i)$。设 $\check X_i$ 是 $X_i$ 如定理 6.2.3 中的超 Kähler 旋转，并设 $\check Y_i$ 是紧化 $\check X_i$ 的有理椭圆曲面。注意 $\check Y_\infty$ 没有定义。

\begin{theorem}{6.6.8}[{Lau--Lee-L., [49]}]
极限 $\lim_{i\to\infty}\check Y_i$ 存在。此外，该极限作为复流形是 $Y_i$ 作为单调辛流形的超势的紧化。\footnote{注意 $Y_i$ 作为单调辛流形的超势不依赖于 $Y_i$ 的复结构。}
\end{theorem}

这一结果或许是 SYZ 纤维化存在性的一个意外推论。回忆在镜像对称的早期，人们曾以为超 Kähler 旋转可能是 Calabi--Yau 曲面镜像对称的底层机制，但这一预期很快被 Gross--Wilson [27] 的工作表明是不正确的。当 $D_t$ 退化到结点曲线时，$X_t$ 在无穷远处的 SYZ 纤维化正在坍缩，于是从 SYZ 观点看 $(Y_0,D_0)$ 可以看作大复结构极限。按这种解释，上述定理大致是说：“超 Kähler 旋转趋向大
% [原文 p.20]
复结构极限的极限给出 del Pezzo 曲面的镜像。”注意这些 $Y_t$ 对不同的 $t$ 都彼此辛同胚，因而具有相同的镜像超势。定理的证明是直接计算两边并作比较：两边都给出形变族中那个特殊的有理椭圆曲面 $\check Y_e$。为这一结果给出直接的 Floer 理论解释将会很有意思。

\subsection{ALH$^*$ 引力瞬子的 Torelli 定理}

引力瞬子是 Hawking 作为其欧氏量子引力理论的基本构件而引入的 [30]。在数学上，引力瞬子就是具有 $L^2$ 曲率张量的非紧完备超 Kähler 四维流形。最早的引力瞬子属于 ALE 型，它们具有最大体积增长。这里 ALE 是“渐近局部欧氏”（asymptotically locally Euclidean）的缩写。后来人们又发现新的引力瞬子类，并按类型加以分类：例如 ALF、ALG、ALH，其体积增长分别为 $r^3$、$r^2$、$r$。第一个是“渐近局部平坦”（asymptotically locally flat）的缩写，后两个则只是按归纳方式命名。后来 Hein [32] 利用有理椭圆曲面构造出超出 ALE/F/G/H 分类的新例子，现在分别称为 ALH$^*$ 与 ALG$^*$。具体地说，Hein 证明：给定任意有理椭圆曲面 $\check Y$ 与纤维 $\check D$，补集 $\check X=\check Y\setminus\check D$ 总容许一个渐近于半平坦度量、具有 $L^2$ 曲率的完全超 Kähler 度量。此外，若 $\check D$ 分别是光滑纤维、具有有限单值的奇异纤维、$I_d^*$ 型奇异纤维或 $I_d$ 型奇异纤维，则该引力瞬子 $\check X$ 相应地为 ALH、ALG、ALG$^*$ 或 ALH$^*$。与我们讨论相关的 ALH$^*$ 引力瞬子具有体积增长 $r^{4/3}$。尽管 ALG$^*$ 引力瞬子的体积增长为 $r^2$，它们的曲率衰减与 ALG 引力瞬子不同。当引力瞬子不具有最大体积增长时，Sun--Zhang [60] 发现其无穷远处由某个 Gibbons--Hawking ansatz 建模，从而断定所有引力瞬子都属于上述六类之一。

人们可能想进一步对每一类中的引力瞬子作分类。这一问题分为三部分：
\begin{enumerate}
\item 给定类型的引力瞬子的微分同胚型有哪些？
\item 在每个微分同胚型之内，给定 $H^2$ 中三个上同调类，它们能否实现为给定类型的引力瞬子的超 Kähler 三元组的上同调类？
\item 若同一微分同胚型、且超 Kähler 三元组上同调类相同的两个引力瞬子，是否必为同一个？
\end{enumerate}
所有能实现为引力瞬子的第二上同调类三元组之集称为\emph{周期区域}。第三个问题问的是：周期区域中的每一点是否恰好参数化一个
% [原文 p.21]
引力瞬子；若成立，通常就称为引力瞬子的 Torelli 定理。

Kronheimer 首先完全分类了 ALE 引力瞬子 [41]：每个 ALE 引力瞬子在超 Kähler 旋转下都双全纯于 $\C^2$ 对 $\mathrm{SU}(2)$ 的有限子群之商的 crepant 消解。此外，对任意在任一 $(-2)$-类上不同时为零的三个第二上同调类，恰存在一个实现它的 ALE 引力瞬子。在 ALG、ALH、ALH$^*$、ALG$^*$ 的情形，引力瞬子（在适当的超 Kähler 旋转下）可以紧化为有理椭圆曲面 [6, 12, 37, 14]。这些结果回答了上述情形下的第一个问题。ALH、ALF、ALH$^*$、ALG、ALG$^*$ 的 Torelli 型定理已在 [6, 7, 12, 15, 16] 中得到证明。在本节余下的部分，我们将把 ALH$^*$ 引力瞬子的 Torelli 定理作为本文所讨论结果的一个快速应用来加以解释。

首先我们回顾 K3 曲面的 Torelli 定理作为类比：所有 K3 曲面都彼此微分同胚 [40]，$(2,0)$-形式的上同调类决定 K3 曲面的复结构 [59, 5, 50]，而 Yau 关于 Calabi 猜想的著名定理 [64] 在给定的 Kähler 类中唯一地决定 Ricci 平坦度量。这些结果依次回答了上面几个问题。我们将遵循同样的策略来证明引力瞬子的 Torelli 定理。

\begin{definition}{6.6.9}
引力瞬子 $(X,\omega,\Omega)$ 称为 ALH$^*$ 型，如果存在微分同胚 $F:X_C\to X$ 使得
\begin{equation}
\|\nabla^k(\omega-(F^{-1})^*\omega_C)\|=O(r^{-k-\varepsilon}),\quad \|\nabla^k(\Omega-(F^{-1})^*\Omega_C)\|=O(r^{-k-\varepsilon}), \tag{6.6.3}
\end{equation}
对某个 $\varepsilon>0$ 成立，其中 $X_C,\omega_C,\Omega_C$ 由第 6.2.1 节定义。
\end{definition}
% [Tutor Note: (6.6.3) 中两个范数都写成了不带下标的 $\|\cdot\|$（页 21 可核对），且第一个式子右端是 $O(r^{-k-\varepsilon})$、第二个也是同样的衰减阶；两式的差别只在被估计的对象 $\omega-(F^{-1})^*\omega_C$ 与 $\Omega-(F^{-1})^*\Omega_C$。]

估计 (6.6.3) 被 Sun--Zhang [60] 改进为指数衰减。因此，与第 6.2 节相同的论证给出如下单值化结果：

\begin{theorem}{6.6.10}
给定 ALH$^*$ 引力瞬子 $(X,\omega,\Omega)$，由 (6.2.2) 给出的超 Kähler 旋转 $(\check X,\check\omega,\check\Omega)$ 可以通过添加一个 $I_d$-纤维而紧化为有理椭圆曲面 $\check Y$。特别地，ALH$^*$ 引力瞬子恰有 10 个微分同胚型。
\end{theorem}

由局部模型计算可知，Kähler 形式 $\check\omega$ 渐近于某个（可能非标准的）半平坦度量。

现在假设有两个 ALH$^*$ 引力瞬子 $(X_i,\omega_i,\Omega_i)$，并有一个微分同胚 $F:X_2\cong X_1$ 认同超 Kähler 三元组的上同调类。由定理 6.6.10，适当的超 Kähler 旋转 $\check X_i$ 可以通过添加 $I_d$-纤维 $\check D_i$ 而紧化为有理椭圆曲面 $\check Y_i$。利用 Friedman--Morgan [20] 的一个引理并计算 Calabi 模型的映射类群，微分同胚 $F:X_2\to X_1$ 可以同伦地延拓为 $\check Y_2\to\check Y_1$，并且 $F$ 把纤维/截面的同调类送到纤维/截面的同调类。
